\section{Aggregation Across Prompts}\label{app:aggregation}

This appendix concerns aggregation across prompts.
Aggregation \emph{within} a single prompt, across permutations of the response ordering, is described in Section~\ref{sec:ordering} and given by Eq.~\eqref{eq:akbar-mor}: the per-byte $\bar{a}_k$ is a mean of per-permutation per-byte rates, not a ratio of permutation-averaged bits to permutation-averaged byte counts.\footnote{An earlier revision of our implementation computed the ratio-of-means form, in which the bits and byte counts at each position are each averaged across permutations before dividing.
After enough permutations, the byte counts at each position approach the mean response length, so that estimator collapses into the total-bits curve rescaled by a single constant, losing the per-permutation per-byte signal that Eq.~\eqref{eq:akbar-mor} preserves.
Switching the implementation to the form in Eq.~\eqref{eq:akbar-mor} shifts the headline numbers in this paper by 0--17\% (in mean-of-ratios' favor), with no qualitative change to any ranking we report.}

The diversity score $D_{Ca_n}$ is defined relative to a single prompt $p$.
To summarize a policy's diversity across a prompt suite $P = \{p_1, \ldots, p_M\}$, one can average:
\begin{equation}
    \bar{D}_{Ca_n}(\pi) = \frac{1}{M}\sum_{j=1}^{M} D_{Ca_n}(\pi, p_j)
\end{equation}

To compare two policies $\pi_1$ and $\pi_2$, the natural scalar is the difference:
\begin{equation}
    \Delta D_{Ca_n} = \bar{D}_{Ca_n}(\pi_1) - \bar{D}_{Ca_n}(\pi_2)
\end{equation}

$\Delta D_{Ca_n}$ answers ``how much more plausibility-weighted residual diversity does $\pi_1$ have than $\pi_2$, on average across prompts?''
No division is involved, so $\Delta D_{Ca_n}$ is stable even when either policy's score is near zero.